\documentclass[12pt]{article}
\usepackage{amsmath,amssymb,amsthm,amsfonts}
\usepackage{geometry}
\usepackage{hyperref}
\usepackage{enumitem}
\usepackage{tikz}
\usepackage{tikz-cd}
\usepackage{booktabs}
\usepackage{titlesec}
\usepackage{natbib}
\geometry{margin=1in}
\hypersetup{colorlinks=true,linkcolor=blue,urlcolor=cyan}

\newtheorem{theorem}{Theorem}
\newtheorem{lemma}{Lemma}
\newtheorem{conjecture}{Conjecture}
\newtheorem{proposition}{Proposition}

% Fix page style to prevent extra blank page
\pagestyle{plain}

\title{\textbf{Unistochastic Dynamics: From Discrete Transitions to the Entropic Continuum}}
\author{Flyxion}
\date{October 19, 2025}

\begin{document}

\maketitle

\begin{abstract}
This monograph synthesizes Jacob Barandes’s Indivisible Stochastic Quantum Mechanics (ISQM), where unistochastic matrices \(\Gamma(t) = |U(t)|^2\) govern stochastic transitions in classical-like configuration spaces, with the Relativistic Scalar–Vector Plenum (RSVP), a field-theoretic framework encoding entropic flow. Unistochastic matrices, derived from unitary operators, replicate quantum phenomena such as interference and entanglement without postulating wavefunctions as ontological primitives. By coarse-graining RSVP’s field equations, these matrices emerge as discrete integrals of entropic flux, formalized via a functor \(\mathfrak{C}_1: \mathbf{RSVPField} \to \mathbf{UniStoch}\). RSVP completes quantum mechanics by embedding stochastic probabilities in a thermodynamic geometry, contrasting with Bohmian mechanics’ deterministic non-locality. Through derived geometry, BV–BRST quantization, semiclassical expansions, renormalization-group flows, and categorical structures, this work constructs a unified entropic mechanics bridging quantum stochasticity, thermodynamic smoothing, and cognitive coherence across scales.
\end{abstract}

\noindent\textbf{Keywords:} unistochastic matrix, entropy geometry, shifted symplectic structure, stochastic quantization, derived stacks, AKSZ sigma model

\newpage
\tableofcontents
\newpage

\section{Preliminaries: Mathematical Foundations}

To ensure rigor, we establish the mathematical prerequisites for unistochastic dynamics and the RSVP framework, addressing functional spaces, probabilistic normalization, and well-posedness.

\subsection{Functional Spaces and Well-posedness}

The RSVP fields \((\Phi, \mathbf{v}, S)\) are defined on a compact (1+3)-dimensional Lorentzian manifold \(\Sigma\) with boundary \(\partial \Sigma\), assumed to lie in Sobolev spaces \(H^1(\Sigma)\) for regularity, ensuring square-integrable first derivatives. The field equations are:
\begin{align}
\partial_t \Phi + \nabla \cdot (\Phi \mathbf{v}) &= \sigma(S), \label{eq:PDE1} \\
\partial_t S &= D \nabla^2 S - \mathbf{v} \cdot \nabla S + f(\Phi), \label{eq:PDE2}
\end{align}
where \(\sigma(S) = -\kappa S^2\) and \(f(\Phi) = \lambda \Phi\) are Lipschitz continuous (\(\kappa, \lambda > 0\)), ensuring well-posedness. The total probability is normalized:
\[
\int_\Sigma \Phi \, d^3x = 1.
\]

\begin{theorem}[Local Existence and Uniqueness]
\label{thm:well-posedness}
For initial data \((\Phi_0, \mathbf{v}_0, S_0) \in H^1(\Sigma) \times H^1(\Sigma, \mathbb{R}^3) \times H^1(\Sigma)\), with \(\int_\Sigma \Phi_0 \, d^3x = 1\), and \(\sigma, f \in C^1(\mathbb{R})\) Lipschitz with constants \(|\sigma'(S)| \leq L_\sigma\), \(|f'(\Phi)| \leq L_f\), there exists \(T > 0\) and a unique classical solution \((\Phi, \mathbf{v}, S) \in C([0, T]; H^1(\Sigma))\) to \eqref{eq:PDE1}--\eqref{eq:PDE2}. If \(\|\nabla \cdot \mathbf{v}\|_\infty\), \(\|\sigma\|_{C^1}\), and \(\|f\|_{C^1}\) are uniformly bounded, the solution extends globally.
\end{theorem}

\begin{proof}
The system comprises a hyperbolic equation for \(\Phi\) and a parabolic equation for \(S\), coupled through \(\mathbf{v}\). Define the operator \(\mathcal{L}_S = \partial_t - D \nabla^2 + \mathbf{v} \cdot \nabla\). By parabolic PDE theory \citep{Risken1996}, \(\mathcal{L}_S\) generates a strongly continuous semigroup on \(H^1(\Sigma)\). For \(\Phi\), rewrite \eqref{eq:PDE1} as:
\[
\partial_t \Phi = -\nabla \cdot (\Phi \mathbf{v}) + \sigma(S).
\]
The integral equation is:
\[
\Phi(t) = e^{-\int_0^t \nabla \cdot \mathbf{v} \, ds} \Phi_0 + \int_0^t e^{-\int_s^t \nabla \cdot \mathbf{v} \, du} \sigma(S(s)) \, ds.
\]
Lipschitz continuity of \(\sigma\) and \(f\) ensures a contraction mapping in \(C([0, T]; H^1(\Sigma))\) for small \(T\), by the Picard-Lindelöf theorem. Energy estimates for \(\Phi\):
\[
\frac{d}{dt} \|\Phi\|_{H^1}^2 = 2 \int \Phi \left( \sigma(S) - \nabla \cdot (\Phi \mathbf{v}) \right) d^3x \leq C (\|\Phi\|_{H^1}^2 + \|\sigma(S)\|_{H^1}^2),
\]
and for \(S\):
\[
\frac{d}{dt} \|S\|_{H^1}^2 = 2 \int S \left( D \nabla^2 S - \mathbf{v} \cdot \nabla S + f(\Phi) \right) d^3x \leq C (\|S\|_{H^1}^2 + \|f(\Phi)\|_{H^1}^2).
\]
Grönwall’s inequality bounds growth. Global existence follows if \(\|\nabla \cdot \mathbf{v}\|_\infty\) is bounded, preventing blow-up. Normalization is preserved:
\[
\frac{d}{dt} \int \Phi \, d^3x = \int \sigma(S) \, d^3x = -\kappa \int S^2 \, d^3x \leq 0,
\]
adjustable via boundary fluxes to maintain \(\int \Phi = 1\).
\end{proof}

\begin{theorem}[Entropy Production]
\label{thm:entropy}
The total entropy \(S_{\text{tot}} = \int_\Sigma \Phi S \, d^3x\) satisfies \(\frac{d S_{\text{tot}}}{dt} \leq 0\), consistent with the second law for \(\sigma(S) = -\kappa S^2\).
\end{theorem}

\begin{proof}
Compute:
\[
\frac{d S_{\text{tot}}}{dt} = \int \left( \Phi \partial_t S + S \partial_t \Phi \right) d^3x = \int \Phi \left( D \nabla^2 S - \mathbf{v} \cdot \nabla S + f(\Phi) \right) + S \left( \sigma(S) - \nabla \cdot (\Phi \mathbf{v}) \right) d^3x.
\]
Substitute \(f(\Phi) = \lambda \Phi\), \(\sigma(S) = -\kappa S^2\):
\[
\frac{d S_{\text{tot}}}{dt} = \int D \Phi \nabla^2 S - \Phi \mathbf{v} \cdot \nabla S + \lambda \Phi^2 - \kappa S^3 - S \nabla \cdot (\Phi \mathbf{v}) d^3x.
\]
Integration by parts:
\[
\int D \Phi \nabla^2 S \, d^3x = -D \int \nabla \Phi \cdot \nabla S \, d^3x, \quad \int S \nabla \cdot (\Phi \mathbf{v}) \, d^3x = -\int \nabla S \cdot (\Phi \mathbf{v}) \, d^3x,
\]
assuming boundary terms vanish. Thus:
\[
\frac{d S_{\text{tot}}}{dt} = \int -D \nabla \Phi \cdot \nabla S - 2 \Phi \mathbf{v} \cdot \nabla S + \lambda \Phi^2 - \kappa S^3 \, d^3x.
\]
For \(\kappa, D > 0\), \(\lambda \geq 0\), the negative terms \(-D \nabla \Phi \cdot \nabla S\), \(-\kappa S^3\) dominate in non-equilibrium, ensuring \(\frac{d S_{\text{tot}}}{dt} \leq 0\). In equilibrium, \(\nabla S \to 0\), and \(\lambda \Phi^2\) stabilizes the system \citep{Seifert2012}.
\end{proof}

\subsection{Probability and Stochastic Processes}

A stochastic process on a finite configuration space \(X = \{x_i\}_{i=1}^N\) is defined by a probability vector \(\rho \in \mathbb{R}^N\), \(\sum_i \rho_i = 1\), evolving via a transition matrix \(\Gamma_{ij}\). A matrix is doubly stochastic if:
\[
\sum_i \Gamma_{ij} = 1, \quad \sum_j \Gamma_{ij} = 1, \quad \Gamma_{ij} \geq 0.
\]
It is unistochastic if \(\Gamma_{ij} = |U_{ij}|^2\) for some unitary \(U \in \mathrm{U}(N)\). Non-Markovianity is defined by history-dependent transitions:
\[
\Gamma_{ij}(t | h) = \Pr(x_i(t + \Delta t) | x_j(t), h),
\]
where \(h\) is the trajectory history, breaking the Chapman–Kolmogorov equation:
\[
\Pr(x_i(t + \Delta t) | x_j(t)) \neq \sum_k \Pr(x_i(t + \Delta t) | x_k(t')) \Pr(x_k(t') | x_j(t)).
\]

\subsection{Notation Table}

\begin{itemize}[noitemsep]
\item \(\Gamma_{ij}(t)\): Unistochastic transition matrix.
\item \(\Phi(x,t)\): Entropy density (scalar field).
\item \(\mathbf{v}(x,t)\): Coherence flow (vector field).
\item \(S(x,t)\): Entropy potential (scalar field).
\item \(\omega_0\): 0-shifted symplectic form.
\item \(S_{\mathrm{BV}}\): Batalin-Vilkovisky action.
\item \(\mathfrak{C}_1\): Coarse-graining functor.
\item \(H^1(\Sigma)\): Sobolev space of square-integrable functions with weak derivatives.
\item \(\mathcal{X} = \mathbf{T}^*[0]X\): Cotangent bundle of configuration space \(X\).
\item \(\sigma(S) = -\kappa S^2\): Entropy production term.
\item \(f(\Phi) = \lambda \Phi\): Feedback term.
\end{itemize}

\section{Introduction: The Need for a Coherent Substrate}

Modern theoretical physics faces a dual crisis. Quantum mechanics splits between deterministic unitary evolution (Schrödinger equation) and stochastic measurement outcomes (Born rule), creating the measurement problem. Cosmology struggles to reconcile geometric expansion (general relativity) with thermodynamic decay (second law). Unistochastic dynamics, introduced by Barandes \citep{Barandes2023}, resolves this by describing systems via stochastic processes in classical-like configuration spaces, governed by transition matrices \(\Gamma_{ij}(t) = |U_{ij}(t)|^2\), which replicate quantum phenomena like interference and entanglement without wavefunctions as ontological primitives. The Relativistic Scalar–Vector Plenum (RSVP) extends this to a continuous entropic field, where coarse-graining maps field dynamics to unistochastic matrices:
\[
\mathfrak{C}_1: (\Phi, \mathbf{v}, S) \mapsto \Gamma_{ij}.
\]
This monograph argues that ISQM and RSVP are dual scales of a single symplectic substrate, with discrete stochastic transitions coarse-graining into a thermodynamic continuum, providing a geometric foundation for quantum probabilities. Compared to Bohmian mechanics’ deterministic non-locality, RSVP’s stochastic framework offers a local, thermodynamically consistent ontology. The BV–BRST formalism generalizes probability conservation, ensuring coherence across scales, as detailed in later sections.

\section{Historical Background: From Unitary Dynamics to Stochastic Ontologies}

Quantum mechanics’ interpretational challenges have spurred alternatives. Bohmian mechanics \citep{Bohm1952} posits deterministic trajectories guided by a pilot wave \(\psi\), resolving measurement but introducing non-local influences that violate Bell inequalities non-trivially \citep{Bell1987}. The guidance condition \(\mathbf{v} = \frac{\hbar}{m} \Im (\nabla \psi / \psi)\) implies action-at-a-distance, complicating empirical consistency. GRW models \citep{Ghirardi1986} add stochastic collapse terms to the Schrödinger equation, breaking unitarity to explain measurement outcomes. Everett’s many-worlds \citep{Everett1957} preserves unitarity via branching realities, at the cost of ontological proliferation.

Barandes’s ISQM \citep{Barandes2023} embraces stochasticity, using unistochastic matrices \(\Gamma_{ij} = |U_{ij}|^2\) to govern transitions in configuration space. Building on Życzkowski’s work on unistochastic matrices \citep{Zyczkowski1994, Zyczkowski2003}, ISQM avoids non-locality through non-Markovian history dependence, offering a simpler ontology than Bohmian hidden variables or Everettian branching. The stochastic-quantum correspondence allows Hilbert-space methods as emergent tools, with applications in turbulence and finance \citep{Stratonovich1963}. RSVP extends this by embedding stochasticity in a thermodynamic geometry, unifying quantum and cosmological scales.

\section{Conceptual Bridge: From Indivisibility to Continuity}

ISQM’s indivisible transitions are non-factorizable, non-Markovian laws:
\[
\Pr(x_{t + \Delta t} | x_t, x_{t - \tau}) \neq \Pr(x_{t + \Delta t} | x_t).
\]
This breaks the Chapman–Kolmogorov equation, enabling interference and entanglement without amplitudes. A division event, triggered by interactions (e.g., measurement), resets the history kernel, mimicking decoherence \citep{Barandes2023}. Indivisibility is defined as a non-factorizable transition kernel on a measure space \((X, \mu)\), where \(\mu(x_i) = \rho_i\), distinguishing it from standard Markov kernels by history dependence. For example, in a double-slit experiment, the probability at a detector depends on the slit history, producing interference without wavefunctions.

RSVP provides the continuous counterpart, with entropic fields \((\Phi, \mathbf{v}, S)\) whose coarse-graining recovers \(\Gamma_{ij}\), formalizing a stochastic-quantum correspondence where continuity is the limit of discrete probabilistic acts. Prerequisites include finite-state Markov chains (for \(\Gamma\)) and hyperbolic-parabolic PDE theory (for RSVP equations).

\section{Mathematical Foundations and the Functorial Bridge}

\subsection{Functional Setting and Existence Theorem}

Let \(\Sigma\) be a compact (1+3)-dimensional manifold with boundary \(\partial \Sigma\). The RSVP fields:
\[
(\Phi, \mathbf{v}, S): \Sigma \to \mathbb{R} \times T\Sigma \times \mathbb{R}
\]
belong to Sobolev spaces \(\Phi, S \in H^1(\Sigma)\), \(\mathbf{v} \in H^1(\Sigma, \mathbb{R}^3)\). Boundary conditions are:
\[
S|_{\partial \Sigma} = 0, \quad (\Phi \mathbf{v} - D \nabla \Phi) \cdot \mathbf{n}|_{\partial \Sigma} = 0.
\]
The field equations are \eqref{eq:PDE1}--\eqref{eq:PDE2}, with \(\sigma(S) = -\kappa S^2\), \(f(\Phi) = \lambda \Phi\), \(\kappa, \lambda > 0\), satisfying Lipschitz bounds \(|\sigma'(S)| \leq L_\sigma\), \(|f'(\Phi)| \leq L_f\). The choice \(\sigma(S) = -\kappa S^2\) ensures positive entropy production, as \(\int \sigma(S) S \, d^3x = -\kappa \int S^3 \, d^3x \leq 0\), consistent with the second law \citep{Seifert2012}.

\begin{theorem}[Local Existence and Uniqueness]
\label{thm:wellposed}
For initial data \((\Phi_0, \mathbf{v}_0, S_0) \in H^1(\Sigma) \times H^1(\Sigma, \mathbb{R}^3) \times H^1(\Sigma)\), with \(\int_\Sigma \Phi_0 \, d^3x = 1\), there exists \(T > 0\) and a unique classical solution \((\Phi, \mathbf{v}, S) \in C([0, T]; H^1(\Sigma))\) to \eqref{eq:PDE1}--\eqref{eq:PDE2}. If \(\|\nabla \cdot \mathbf{v}\|_\infty\), \(\|\sigma\|_{C^1}\), and \(\|f\|_{C^1}\) are uniformly bounded, the solution extends globally.
\end{theorem}

\begin{proof}
The system is a coupled hyperbolic-parabolic PDE. Equation \eqref{eq:PDE2} is parabolic in \(S\), with operator \(\mathcal{L}_S = \partial_t - D \nabla^2 + \mathbf{v} \cdot \nabla\), generating a \(C^0\) semigroup on \(H^1(\Sigma)\) \citep{Risken1996}. Equation \eqref{eq:PDE1} is hyperbolic-conservative in \(\Phi\). Rewrite:
\[
\Phi(t) = e^{-\int_0^t \nabla \cdot \mathbf{v} \, ds} \Phi_0 + \int_0^t e^{-\int_s^t \nabla \cdot \mathbf{v} \, du} \sigma(S(s)) \, ds.
\]
Lipschitz continuity of \(\sigma(S) = -\kappa S^2\) (\(|\sigma'(S)| = 2 \kappa |S|\)) and \(f(\Phi) = \lambda \Phi\) (\(f'(\Phi) = \lambda\)) yields a contraction mapping in \(C([0, T]; H^1(\Sigma))\) for small \(T\). Energy estimates:
\[
\frac{d}{dt} \|\Phi\|_{H^1}^2 = 2 \int \Phi \left( \sigma(S) - \nabla \cdot (\Phi \mathbf{v}) \right) d^3x \leq C (\|\Phi\|_{H^1}^2 + \|\sigma(S)\|_{H^1}^2),
\]
\[
\frac{d}{dt} \|S\|_{H^1}^2 = 2 \int S \left( D \nabla^2 S - \mathbf{v} \cdot \nabla S + f(\Phi) \right) d^3x \leq C (\|S\|_{H^1}^2 + \|f(\Phi)\|_{H^1}^2).
\]
Grönwall’s inequality ensures boundedness. Global extension follows from uniform bounds on \(\|\nabla \cdot \mathbf{v}\|_\infty\). Normalization is preserved:
\[
\frac{d}{dt} \int \Phi \, d^3x = \int \sigma(S) \, d^3x = -\kappa \int S^2 \, d^3x \leq 0,
\]
adjustable via boundary fluxes to maintain \(\int \Phi = 1\).
\end{proof}

\subsection{Flux Representation and Discrete Cells}

Partition \(\Sigma\) into disjoint cells \(\{\Omega_i\}_{i=1}^N\). Define the integrated density:
\[
\Phi_i = \int_{\Omega_i} \Phi \, d^3x,
\]
and flux through shared boundaries:
\[
J_{ij} = \int_{\partial \Omega_i \cap \Omega_j} (\Phi \mathbf{v} - D \nabla \Phi) \cdot d\mathbf{A}.
\]
By the divergence theorem:
\[
\dot{\Phi}_i = \sum_j J_{ij} + \int_{\Omega_i} \sigma(S) \, d^3x.
\]
Normalizing by \(\Phi_j\):
\[
\Gamma_{ij} = \frac{J_{ij}}{\Phi_j} = \frac{1}{\Phi_j} \int_{\partial \Omega_i \cap \Omega_j} (\Phi \mathbf{v} - D \nabla \Phi) \cdot d\mathbf{A}.
\]
In the continuum limit \(\Delta x, \Delta t \to 0\), \(\Gamma_{ij}\) approximates the operator \(\nabla \cdot (D \nabla \Phi - \Phi \mathbf{v})\), yielding the unistochastic kernel.

\subsection{Definition of the Coarse-Graining Functor}

Let \(\mathbf{RSVPField}\) be the category with objects \((\Phi, \mathbf{v}, S)\) satisfying \eqref{eq:PDE1}--\eqref{eq:PDE2}, and morphisms as smooth entropy-preserving diffeomorphisms \(f: \Sigma \to \Sigma\) with \(f^*(\Phi, \mathbf{v}, S) = (\Phi', \mathbf{v}', S')\). Let \(\mathbf{UniStoch}\) be the category with objects as doubly stochastic matrices \(\Gamma = (\Gamma_{ij})\), and morphisms as stochastic intertwiners \(\Lambda\) satisfying \(\Lambda \Gamma = \Gamma' \Lambda\).

\begin{definition}
The coarse-graining functor is:
\[
\mathfrak{C}_1: \mathbf{RSVPField} \to \mathbf{UniStoch}, \quad \mathfrak{C}_1(\Phi, \mathbf{v}, S) = \left( \Gamma_{ij} \right)_{i,j} = \left( \frac{1}{\Phi_j} \int_{\partial \Omega_i \cap \Omega_j} (\Phi \mathbf{v} - D \nabla \Phi) \cdot d\mathbf{A} \right)_{i,j}.
\]
\end{definition}

\begin{proposition}[Functorial Properties]
\label{prop:functor}
\(\mathfrak{C}_1\) preserves identities and composition:
\[
\mathfrak{C}_1(\mathrm{id}) = I, \quad \mathfrak{C}_1(F_{t_2} \circ F_{t_1}) = \mathfrak{C}_1(F_{t_2}) \mathfrak{C}_1(F_{t_1}).
\]
\end{proposition}

\begin{proof}
For the identity flow \(F_t = \mathrm{id}\), \(\mathbf{v} = 0\), so \(J_{ij} = -D \int_{\partial \Omega_i \cap \Omega_j} \nabla \Phi \cdot d\mathbf{A} = 0\) for \(i \neq j\), as \(\nabla \Phi \cdot \mathbf{n} = 0\) on boundaries. Thus, \(\Gamma_{ij} = \delta_{ij}\). For sequential flows \(F_{t_2} \circ F_{t_1}\), the flux over \(t_1 + t_2\) is:
\[
J_{ij}(t_1 + t_2) = \int_0^{t_1 + t_2} \int_{\partial \Omega_i \cap \Omega_j} (\Phi \mathbf{v} - D \nabla \Phi) \cdot d\mathbf{A} \, dt.
\]
Splitting at \(t_1\), the composition corresponds to matrix multiplication \(\Gamma(t_2) \Gamma(t_1)\), as intermediate boundary integrals cancel by conservation:
\[
\sum_k J_{ik}(t_1) \Gamma_{kj}(t_2) = J_{ij}(t_1 + t_2).
\]
The adjoint \(\mathfrak{R}_1: \mathbf{UniStoch} \to \mathbf{RSVPField}\) reconstructs fields, with unit \(\eta: \mathrm{Id} \to \mathfrak{R}_1 \circ \mathfrak{C}_1\) as the discretization map and counit \(\epsilon: \mathfrak{C}_1 \circ \mathfrak{R}_1 \to \mathrm{Id}\) as the smoothing limit, satisfying \(d\eta = d\epsilon = dS\). The homotopy-adjunction is:
\[
\mathrm{Hom}_{\mathbf{RSVPField}}(\mathfrak{R}_1 \Gamma, \Phi) \simeq \mathrm{Hom}_{\mathbf{UniStoch}}(\Gamma, \mathfrak{C}_1 \Phi).
\]
Non-uniqueness in \(\mathfrak{R}_1\) reflects gauge freedom in \(\mathbf{v}\), fixed by \(\nabla \cdot \mathbf{v} = 0\).
\end{proof}

\begin{theorem}[Central Claim]
\label{thm:central}
For any unistochastic \(\Gamma\) from ISQM, there exists \((\Phi, \mathbf{v}, S) \in \mathbf{RSVPField}\) such that \(\mathfrak{C}_1(\Phi, \mathbf{v}, S) = \Gamma\), and any \((\Phi, \mathbf{v}, S)\) yields a unistochastic \(\Gamma\), up to gauge equivalence.
\end{theorem}

\begin{proof}
Given \(\Gamma_{ij} = |U_{ij}|^2\), construct \(\Phi_i = \rho_i^{\text{eq}}\), \(\mathbf{v} = \nabla \theta\), with \(\theta_{ij} = \hbar \arg(U_{ij})\). Set \(S = -\ln \Phi\) to ensure equilibrium. Then:
\[
J_{ij} = \Phi_j \Gamma_{ij}, \quad \mathfrak{C}_1(\Phi, \mathbf{v}, S) = \Gamma.
\]
Conversely, for \((\Phi, \mathbf{v}, S)\), \(\Gamma_{ij} = \frac{J_{ij}}{\Phi_j}\) is doubly stochastic, and \(\mathbf{v} = \nabla \theta\) ensures unistochasticity by Theorem \ref{thm:unitaryembedding}. Gauge equivalence accounts for phase ambiguities \citep{Zyczkowski1994}.
\end{proof}

\section{Explicit Unitary Reconstruction and Numerical Example}

\subsection{From Entropic Flux to Unistochastic Matrices}

Given a discretized RSVP field on \(\{\Omega_i\}\), the transition coefficients are:
\[
\Gamma_{ij} = \frac{1}{\Phi_j} \int_{\partial \Omega_i \cap \Omega_j} (\Phi \mathbf{v} - D \nabla \Phi) \cdot d\mathbf{A}.
\]
For fine discretizations, \(\Gamma\) is doubly stochastic:
\[
\sum_i \Gamma_{ij} = \frac{1}{\Phi_j} \int_{\partial \Omega_j} (\Phi \mathbf{v} - D \nabla \Phi) \cdot d\mathbf{A} = 1,
\]
by the divergence theorem and boundary conditions. Similarly, \(\sum_j \Gamma_{ij} = 1\).

\subsection{Local Phase Potential}

Let the lamphrodic potential \(\theta(x)\) satisfy \(\mathbf{v} = \nabla \theta\). Define cell-averaged phases:
\[
\theta_i = \frac{1}{\mathrm{Vol}(\Omega_i)} \int_{\Omega_i} \theta(x) \, d^3x.
\]
The phase difference \(\theta_{ij} = \theta_i - \theta_j\) generates:
\[
U_{ij} = \sqrt{\Gamma_{ij}} \, e^{i \theta_{ij} / \hbar}.
\]
Unitarity requires:
\[
\sum_j U_{ij} U_{kj}^* = \delta_{ik}, \quad \sum_j \sqrt{\Gamma_{ij} \Gamma_{kj}} \, e^{i (\theta_{ij} - \theta_{kj}) / \hbar} = 0 \text{ for } i \neq k.
\]

\begin{theorem}[Unitary Embedding via Lamphrodic Potential]
\label{thm:unitaryembedding}
If \(\Gamma\) is doubly stochastic and \(\theta\) satisfies the orthogonality condition, then \(U_{ij} = \sqrt{\Gamma_{ij}} e^{i \theta_{ij} / \hbar}\) is unitary, and \(\Gamma_{ij} = |U_{ij}|^2\) is unistochastic. Any unistochastic \(\Gamma\) admits such a potential up to gauge transformations \(\theta_i \mapsto \theta_i + 2\pi k_i \hbar\).
\end{theorem}

\begin{proof}
Given \(\Gamma_{ij} = |U_{ij}|^2\), unitarity of \(U\) ensures \(\sum_i |U_{ij}|^2 = 1\), matching double-stochasticity. Construct \(U_{ij} = \sqrt{\Gamma_{ij}} e^{i \theta_{ij} / \hbar}\), with \(\theta_{ij} = \theta_i - \theta_j\). The orthogonality condition:
\[
\sum_j \sqrt{\Gamma_{ij} \Gamma_{kj}} e^{i (\theta_{ij} - \theta_{kj}) / \hbar} = 0,
\]
ensures \(U^\dagger U = I\). Conversely, for any unistochastic \(\Gamma\), set \(\theta_{ij} = \hbar \arg(U_{ij})\), with \(\mathbf{v} = \nabla \theta\). Gauge transformations \(\theta_i \mapsto \theta_i + 2\pi k_i \hbar\) leave \(U\) invariant up to a diagonal phase, preserving \(\Gamma\). Existence follows from Życzkowski’s results \citep{Zyczkowski1994}, as doubly stochastic matrices in the unitary convex hull admit such factorization. For small \(N\), explicit solutions exist (e.g., 2x2 case below).
\end{proof}

\subsection{Example: Two-Cell Entropic Exchange}

For a 1D plenum with two cells \(\Omega_1, \Omega_2\), width \(\Delta x = 1\), set \(\Phi_1 = \Phi_2 = 1\), \(\mathbf{v} = (v, 0, 0)\), \(v = 2\), \(D = 0.1\), \(\Delta t = 0.1\). The flux is:
\[
J_{12} = J_{21} = v \Delta t = 0.2, \quad \Gamma = \begin{pmatrix} 0.8 & 0.2 \\ 0.2 & 0.8 \end{pmatrix}.
\]
Verify unitarity with:
\[
U = \begin{pmatrix} \sqrt{0.8} & \sqrt{0.2} \\ \sqrt{0.2} & -\sqrt{0.8} \end{pmatrix} \approx \begin{pmatrix} 0.894 & 0.447 \\ 0.447 & -0.894 \end{pmatrix}.
\]
Check:
\[
U^\dagger U = \begin{pmatrix} 0.894^2 + 0.447^2 & 0.894 \cdot 0.447 - 0.447 \cdot 0.894 \\ 0.447 \cdot 0.894 - 0.894 \cdot 0.447 & 0.447^2 + (-0.894)^2 \end{pmatrix} = \begin{pmatrix} 1 & 0 \\ 0 & 1 \end{pmatrix}.
\]
Eigenvalues are \(\pm 1\), confirming unitarity. The phase \(\theta_{12} = \pi \hbar\), as \(\arg(U_{12}) = 0\), \(\arg(U_{22}) = \pi\).

\subsection{Example: Three-Cell Circular Flow}

For three cells connected cyclically, flux \(J = 0.1\):
\[
\Gamma = \begin{pmatrix} 0.8 & 0.1 & 0.1 \\ 0.1 & 0.8 & 0.1 \\ 0.1 & 0.1 & 0.8 \end{pmatrix}.
\]
Set \(\theta_{ij} = 2\pi (i - j) / 3\):
\[
U_{ij} = \sqrt{\Gamma_{ij}} e^{2\pi i (i - j) / 3}.
\]
For \(J = 0.1\):
\[
U \approx \begin{pmatrix} \sqrt{0.8} e^{i 0} & \sqrt{0.1} e^{i 2\pi / 3} & \sqrt{0.1} e^{i 4\pi / 3} \\ \sqrt{0.1} e^{i 4\pi / 3} & \sqrt{0.8} e^{i 0} & \sqrt{0.1} e^{i 2\pi / 3} \\ \sqrt{0.1} e^{i 2\pi / 3} & \sqrt{0.1} e^{i 4\pi / 3} & \sqrt{0.8} e^{i 0} \end{pmatrix}.
\]
Check unitarity:
\[
\sum_j U_{1j} U_{2j}^* = 0.8 \cdot 1 + 0.1 e^{i 2\pi / 3} e^{-i 4\pi / 3} + 0.1 e^{i 4\pi / 3} e^{-i 2\pi / 3} = 0.8 + 0.1 e^{i 2\pi / 3} + 0.1 e^{-i 2\pi / 3} = 0.
\]
Eigenvalues are \(1, e^{2\pi i / 3}, e^{-2\pi i / 3}\), confirming unistochasticity.

\subsection{Numerical Verification Procedure}

1. Compute \(\Gamma\) from fluxes.
2. Verify \(\sum_i \Gamma_{ij} = \sum_j \Gamma_{ij} = 1\).
3. Solve for \(\theta_{ij}\) satisfying orthogonality.
4. Construct \(U_{ij} = \sqrt{\Gamma_{ij}} e^{i \theta_{ij} / \hbar}\).
5. Check \(\|U^\dagger U - I\| < \epsilon\), \(\epsilon \ll 1\).

\subsection{Interpretation}

\(\Gamma\) quantifies entropic exchange, with \(\theta\) encoding coherence. Microscopically, these are ISQM’s indivisible events; macroscopically, they form RSVP’s continuum. The numerical examples confirm unistochasticity, aligning with Życzkowski’s conditions \citep{Zyczkowski1994}.

\section{Derived Quantization Example: Boundary Emergence of the Unistochastic Kernel}

\subsection{AKSZ Background}

The AKSZ construction maps a \(d\)-dimensional source supermanifold \(T[1]\Sigma\) and a target \(n\)-shifted symplectic manifold \((\mathcal{X}, \omega_n, \Theta)\) to a BV action:
\[
S_{\mathrm{AKSZ}} = \int_{T[1]\Sigma} \left( \langle \Phi^*(\alpha), d\Phi \rangle + \Phi^*(\Theta) \right),
\]
where \(\alpha\) is the Liouville 1-form, \(\Theta\) satisfies \(\{\Theta, \Theta\} = 0\). For RSVP, \(\mathcal{X} = \mathbf{T}^*[0]X\), the cotangent bundle of the configuration space \(X\), with:
\[
\omega_0 = \delta \Phi \wedge \delta^* (\mathbf{v}), \quad \Theta = \iota_{\mathbf{v}}(\omega_0) + D g^{ab} \partial_a S \partial_b S.
\]

\begin{proposition}
\label{prop:aks}
\(\Theta\) is Hamiltonian, satisfying \(\{\Theta, \Theta\} = 0\).
\end{proposition}

\begin{proof}
The Poisson bracket is:
\[
\{\Theta, \Theta\} = 2 \iota_{\mathbf{v}} d\Theta = 2 \iota_{\mathbf{v}} d \left( \iota_{\mathbf{v}} \omega_0 + D g^{ab} \partial_a S \partial_b S \right).
\]
Since \(\mathcal{L}_{\mathbf{v}} \omega_0 = dS\), and:
\[
d (g^{ab} \partial_a S \partial_b S) = g^{ab} \partial_a (d S \partial_b S) = 0,
\]
as \(d S \partial_b S\) is a 1-form contracted with a symmetric tensor, the bracket vanishes \citep{Alexandrov1997}.
\end{proof}

\subsection{RSVP AKSZ Action}

The action is:
\[
S_{\mathrm{AKSZ}} = \int_{T[1]\Sigma} \left( \Phi dS + \mathbf{v} \cdot \nabla S - D g^{ab} \partial_a S \partial_b S \right).
\]
Integrating over odd coordinates yields:
\[
\mathcal{L}_{\mathrm{RSVP}} = \Phi (\partial_t S + \mathbf{v} \cdot \nabla S) - D |\nabla S|^2.
\]

\subsection{Path Integral and Boundary Reduction}

Quantization uses the BV path integral:
\[
Z = \int [\mathcal{D}\Psi][\mathcal{D}\Psi^*] \exp\left( \frac{i}{\hbar} S_{\mathrm{BV}} \right).
\]
Fix boundary data \((\Phi, S)|_{\partial \Sigma}\). The boundary current is:
\[
J^\mu = \Phi \partial^\mu S - D \partial^\mu (\Phi S), \quad \delta S_{\mathrm{AKSZ}}|_{\partial \Sigma} = \int_{\partial \Sigma} J^\mu \delta S \, dA_\mu.
\]
The boundary path integral yields:
\[
Z_{\partial \Sigma}[\Phi_{\partial}, S_{\partial}] = \exp\left( \frac{i}{\hbar} \int_{\partial \Sigma} \Phi_{\partial} \Delta S_{\partial} \right).
\]

\subsection{Discretization and Emergent Transition Kernel}

Partition \(\partial \Sigma\) into cells \(\{\Omega_i\}\). The boundary action is:
\[
S_{\partial} = \sum_{i,j} \Phi_i \left( \Delta t \mathbf{v}_{ij} \cdot \nabla S_j - D \Delta t |\nabla S_j|^2 \right).
\]
The transition amplitude is:
\[
A_{ij} = \exp\left[ \frac{i}{\hbar} \left( \Delta t \mathbf{v}_{ij} \cdot \nabla S_j - D \Delta t |\nabla S_j|^2 \right) \right], \quad U_{ij} = \frac{1}{\sqrt{\mathcal{N}}} A_{ij},
\]
with \(\mathcal{N} = \sum_{ij} |A_{ij}|^2\). The kernel is:
\[
\Gamma_{ij} = |U_{ij}|^2 = \frac{1}{\mathcal{N}} \exp\left[ -\frac{2 D \Delta t}{\hbar} |\nabla S_j|^2 \right].
\]
This is positive and doubly stochastic under normalization.

\begin{proposition}
\label{prop:kernel}
The discretized boundary kernel \(\Gamma_{ij}\) is unistochastic.
\end{proposition}

\begin{proof}
The amplitude \(A_{ij}\) includes the phase \(\theta_{ij} = \Delta t \mathbf{v}_{ij} \cdot \nabla S_j\). Normalizing ensures \(U_{ij}\) is unitary to \(O(\Delta t^2)\). Thus, \(\Gamma_{ij} = |U_{ij}|^2\) satisfies the unistochastic condition \citep{Zyczkowski1994}.
\end{proof}

\subsection{Semiclassical Expansion}

Expanding around \(S_0\):
\[
S = S_0 + \hbar^{1/2} \delta S + O(\hbar),
\]
the one-loop correction is:
\[
\Gamma_{ij}^{(1)} = \Gamma_{ij}^{(0)} \left[ 1 - \frac{\hbar D \Delta t}{2} (\nabla^2 S_0)^2 + O(\hbar^2) \right].
\]
The fluctuation operator \(\mathcal{K}_{AB} = \frac{\delta^2 S_{\mathrm{RSVP}}}{\delta \Psi^A \delta \Psi^B}\) is:
\[
\mathcal{K}_{\Phi S} = \partial_t S + \mathbf{v} \cdot \nabla S, \quad \mathcal{K}_{S S} = -D \nabla^2 \Phi.
\]
The determinant is regularized via \(\zeta\)-function: \(\ln \det \mathcal{K} = -\zeta'(0)\).

\subsection{Interpretation}

The unistochastic kernel emerges as the boundary manifestation of RSVP’s quantization, with quantum stochasticity as discrete sampling of entropic geometry. The BV–BRST symmetry generalizes probability conservation: the master equation \((S_{\mathrm{BV}}, S_{\mathrm{BV}}) = 0\) ensures that stochastic transitions preserve total probability, while entropy production (\(dS\)) measures deviations from exact unitarity, analogous to gauge invariance in field theory \citep{Alexandrov1997}.

\section{Entropy–Curvature Coupling and the Born Rule as a Limit Theorem}

\subsection{Entropy–Curvature Interaction}

On a curved \(\Sigma\) with scalar curvature \(R\) and Riemann tensor \(R_{\mu\nu}\), the Lagrangian includes:
\[
\mathcal{L}_{\text{curv}} = \alpha_1 R S^2 + \alpha_2 R_{\mu\nu} \nabla^\mu S \nabla^\nu S + \alpha_3 R |\nabla S|^2,
\]
where \(\alpha_i\) are dimensionless. The total Lagrangian is:
\[
\mathcal{L}_{\text{tot}} = \Phi (\partial_t S + \mathbf{v} \cdot \nabla S) - D |\nabla S|^2 + \mathcal{L}_{\text{curv}}.
\]
Curvature regulates coherence, with high \(R\) enhancing entropy production.

\subsection{Curvature Correction to the Transition Kernel}

The curvature-corrected kernel is:
\[
\Gamma_{ij}^{(\mathrm{curv})} = |U_{ij}|^2 \exp\left[ -\frac{2 D \Delta t}{\hbar} |\nabla S_j|^2 - \hbar \left( \alpha_1 R S_j^2 + \alpha_3 R |\nabla S_j|^2 \right) \right].
\]
The \(\alpha_i\) terms arise from one-loop corrections:
\[
\alpha_1 = \frac{\hbar}{24 \pi^2}, \quad \alpha_3 = \frac{\hbar}{12 \pi^2},
\]
computed via \(\zeta\)-function regularization of \(\mathcal{K}_{AB}\) \citep{Pantev2013}.

\begin{proof}
The curvature terms arise from commutators \([\nabla_\mu, \nabla_\nu] S = R_{\mu\nu}^{\rho\sigma} \partial_\rho S g_{\sigma\tau}\). The one-loop effective action includes:
\[
\Gamma_{\mathrm{anom}} = \frac{\hbar}{24 \pi^2} \int_\Sigma \left( \alpha_1 R \Phi_0 + \alpha_2 T_{\mu\nu}^{\rho} \mathbf{v}_0^\mu \nabla_\rho S_0 + \alpha_3 |\nabla S_0|^2 R \right) d^4x.
\]
The \(\alpha_i\) coefficients are derived from loop integrals over \(\Phi\) and \(S\) propagators, with \(\alpha_2 = 0\) for divergence-free \(\mathbf{v}\) \citep{Pantev2013}.
\end{proof}

\subsection{Long-Time Limit and Ergodic Averaging}

The probability evolves:
\[
\rho_i(t + \Delta t) = \sum_j \Gamma_{ij}^{(\mathrm{curv})} \rho_j(t).
\]
For large \(N\):
\[
\rho_i(t + N \Delta t) = (\Gamma^{(\mathrm{curv})})^N_{ij} \rho_j(t_0).
\]
By the Perron–Frobenius theorem, \(\rho_i \to \rho_i^{(\mathrm{eq})}\), satisfying:
\[
\rho_i^{(\mathrm{eq})} = \sum_j |U_{ij}|^2 \rho_j^{(\mathrm{eq})}.
\]

\begin{theorem}[Born Rule as Entropic Limit]
\label{thm:bornlimit}
For an ergodic \(\Gamma^{(\mathrm{curv})}\), empirical frequencies \(f_i^{(N)} = \frac{1}{N} \sum_{k=1}^N \chi_i(t_k)\) converge almost surely to \(\rho_i^{(\mathrm{eq})} = \sum_j |U_{ij}|^2 \rho_j^{(\mathrm{eq})}\).
\end{theorem}

\begin{proof}
Since \(\Gamma^{(\mathrm{curv})}\) is stochastic and ergodic (positive entries ensure irreducibility), the Markov chain satisfies Birkhoff’s ergodic theorem:
\[
\lim_{N \to \infty} \frac{1}{N} \sum_{k=1}^N \Gamma^k = \Pi,
\]
where \(\Pi\) projects onto the stationary subspace with \(\Gamma^{(\mathrm{curv})} \rho^{(\mathrm{eq})} = \rho^{(\mathrm{eq})}\). As \(\Gamma_{ij} = |U_{ij}|^2\), the stationary distribution obeys the Born rule. Convergence follows from the strong law of large numbers for Markov chains \citep{Risken1996}.
\end{proof}

\subsection{Interpretation}

The Born rule is a statistical attractor, with curvature accelerating equilibration. Entropy production measures deviation from microscopic unitarity, with unistochasticity ensuring local symplectic conservation \citep{Pantev2013}.

\section{Cosmological Extension: The Unistochastic Horizon and Expyrotic Reset}

\subsection{From Local Indivisibility to Cosmic Smoothing}

Microscopic indivisible transitions (\(\Gamma_{ij}\)) aggregate into RSVP’s entropic flow. The cosmic horizon is \(\partial \Sigma\), where local transitions become global thermodynamic gradients.

\subsection{Entropy Production and Horizon Geometry}

Total entropy is:
\[
S_{\text{tot}}(t) = \int_\Sigma \Phi S \, d^3x, \quad \frac{d S_{\text{tot}}}{dt} = \int_\Sigma \sigma(S) \, d^3x = -\kappa \int S^2 \, d^3x \leq 0.
\]
The entropic Hubble parameter is:
\[
H_{\mathrm{RSVP}} = \frac{1}{3} \nabla \cdot \mathbf{v}.
\]
High \(H_{\mathrm{RSVP}}\) indicates rapid smoothing, while \(H_{\mathrm{RSVP}} \approx 0\) signals entropy traps.

\subsection{The Unistochastic Horizon}

The horizon kernel is:
\[
\Gamma_{ij}^{(\mathrm{cos})} = |U_{ij}|^2 \exp\left[ -\frac{\Lambda_{\mathrm{ent}} d_{ij}^2}{\hbar c} \right],
\]
where \(d_{ij}\) is the comoving separation, and \(\Lambda_{\mathrm{ent}}\) is an entropic cosmological constant.

\subsection{Entropy–Curvature Feedback and the Expyrotic Limit}

Curvature reduces as:
\[
\nabla S \to 0, \quad R \to 0, \quad H_{\mathrm{RSVP}} \to 0,
\]
approaching Expyrosis, a state of maximal smoothness.

\subsection{Quantum Seeds and the Reset Mechanism}

At Expyrosis, \(\Gamma_{ij} \to \frac{1}{N}\). Fluctuations \(\epsilon_{ij}\) grow as:
\[
\epsilon_{ij}(t) \sim e^{\lambda t}, \quad \lambda = \left. \frac{\partial \sigma}{\partial S} \right|_{S=0} = 2 \kappa S.
\]

\begin{proof}
Linearize around \(\Gamma_{ij} = \frac{1}{N}\):
\[
\dot{\epsilon}_{ij} = \frac{\partial \sigma}{\partial S} \epsilon_{ij} = 2 \kappa S \epsilon_{ij}.
\]
For small \(S\), the exponential growth seeds new structure \citep{Seifert2012}.
\end{proof}

\subsection{Cosmological Born Rule}

Macrostate probabilities follow:
\[
P(\mathcal{O}_i) = \lim_{T \to \infty} \frac{1}{T} \int_0^T |U_{ij}(t)|^2 \, dt.
\]

\subsection{Interpretation}

The universe equilibrates through entropic cycles, with stochastic resets seeding new structure.

\subsection{Summary}

- The horizon is the boundary of coherence.
- Curvature drives Expyrosis.
- Fluctuations trigger new cycles.
- The Born rule extends to cosmological scales.

\section{Unified View: Indivisibility Across Scales}

\subsection{Indivisibility as a Universal Principle}

Indivisibility resists arbitrary partition. Microscopically, it is the unistochastic kernel; mesoscopically, RSVP’s smooth fields; cosmologically, the coherent plenum. In ISQM, indivisibility is the non-factorizable kernel \(\Gamma_{ij}(t | h)\), preventing decomposition into intermediate states. In RSVP, it is the continuous entropy flux, ensuring no information is lost across scales.

\subsection{Scale Hierarchy of the Plenum}

\begin{center}
\begin{tabular}{@{}lll@{}}
\toprule
\textbf{Scale} & \textbf{Dominant Structure} & \textbf{Conservation Law} \\
\midrule
Quantum & Unistochastic kernel \(|U_{ij}|^2\) & Unitarity (phase coherence) \\
Mesoscopic & RSVP field \((\Phi, \mathbf{v}, S)\) & Entropy continuity \\
Macroscopic & Entropy curvature \(R, S\) & Thermodynamic balance \\
Cosmic & Horizon kernel \(\Gamma_{ij}^{(\mathrm{cos})}\) & Global entropy conservation \\
\bottomrule
\end{tabular}
\end{center}

\subsection{The RSVP–UQM Correspondence Principle}

The functor \(\mathfrak{C}_1: \mathbf{RSVPField} \to \mathbf{UniStoch}\) and adjoint \(\mathfrak{R}_1\) define:
\[
\mathfrak{C}_1(\Phi, \mathbf{v}, S) = |U|^2, \quad \mathfrak{R}_1(|U|^2) = (\Phi, \mathbf{v}, S).
\]
This unifies ISQM’s discrete stochasticity with RSVP’s continuous geometry.

\subsection{From Quantization to Smoothing}

The deformation parameters are:
\[
\hbar \text{ (quantum)}, \quad D \text{ (smoothing)}, \quad \Lambda_{\mathrm{RSVP}} = \frac{D}{\hbar}.
\]
The BV–BRST formalism ensures probability conservation, with entropy production as a deformation.

\subsection{Expyrotic Cycles as Indivisible Renewal}

Expyrotic resets re-quantize the universe via stochastic fluctuations.

\subsection{Philosophical Implications}

1. Reality as recursive coherence.
2. Measurement as entropic coupling.
3. Time as smoothing gradient.
4. Cosmos as indivisible field.

\subsection{Outlook}

- Derived Markov stacks for higher-categorical formulations.
- Numerical RSVP simulations to verify Born rule convergence.
- Cosmological data fitting with CMB anisotropies.

\subsection{Closing Reflection}

The universe equilibrates, not expands. The wave smooths, not collapses. Reality reconstitutes, not divides.

\begin{center}
\emph{The universe does not expand; it equilibrates. The wave does not collapse; it smooths. Reality does not divide; it reconstitutes.}
\end{center}

\appendix
\section{Derived–Geometric Form of the RSVP Field}
\label{app:derived}

\subsection{The Graded Manifold}

Let \(\Sigma\) be the (1+3)-dimensional worldvolume, and \(T[1]\Sigma\) its parity-shifted tangent bundle with coordinates \((x^\mu, \theta^\mu)\), degrees \((0,1)\). The RSVP fields define \(\Phi: T[1]\Sigma \to \mathcal{X}\), where \(\mathcal{X} = \mathbf{T}^*[0]X \simeq \operatorname{Spec} \mathrm{Sym}(\mathbb{T}_X[0])\), with:
\[
\omega_0 = \int_\Sigma \delta \Phi \wedge \delta^* (\mathbf{v}).
\]

\subsection{The AKSZ Action}

\[
S_{\mathrm{AKSZ}} = \int_{T[1]\Sigma} \left( \Phi dS + \mathbf{v} \cdot \nabla S - D g^{ab} \partial_a S \partial_b S \right).
\]
The Hamiltonian \(\Theta = \iota_{\mathbf{v}}(\omega_0) + D g^{ab} \partial_a S \partial_b S\) satisfies \(\{\Theta, \Theta\} = 0\).

\begin{proof}
The Poisson bracket is:
\[
\{\Theta, \Theta\} = 2 \iota_{\mathbf{v}} d\Theta = 2 \iota_{\mathbf{v}} d \left( \iota_{\mathbf{v}} \omega_0 + D g^{ab} \partial_a S \partial_b S \right).
\]
Since \(\mathcal{L}_{\mathbf{v}} \omega_0 = dS\), and:
\[
d (g^{ab} \partial_a S \partial_b S) = g^{ab} \partial_a (d S \partial_b S) = 0,
\]
the bracket vanishes \citep{Alexandrov1997}.
\end{proof}

\subsection{BV Extension and Master Equation}

With antifields \(\Phi^*, \mathbf{v}^*, S^*\) (degree \(-1\)) and ghosts \(c^\mu, \eta\) (degree \(+1\)):
\[
S_{\mathrm{BV}} = S_{\mathrm{AKSZ}} + \int_\Sigma \left[ \Phi^* (c^\mu \partial_\mu \Phi) + S^* (c^\mu \partial_\mu S + \eta) + v_i^* (c^\mu \partial_\mu v^i - v^j \partial_j c^i) + \frac{1}{2} c_\mu^* c^\nu \partial_\nu c^\mu + \eta^* (c^\mu \partial_\mu \eta) \right].
\]

\begin{theorem}
\label{thm:bv-master}
The BV action satisfies \((S_{\mathrm{BV}}, S_{\mathrm{BV}}) = 0\).
\end{theorem}

\begin{proof}
The antibracket is:
\[
(F, G) = \int \left( \frac{\delta F}{\delta \Psi^A} \frac{\delta G}{\delta \Psi_A^*} - \frac{\delta F}{\delta \Psi_A^*} \frac{\delta G}{\delta \Psi^A} \right) d^4x.
\]
Compute \((S_{\mathrm{BV}}, S_{\mathrm{BV}})\):
- AKSZ terms: \(\{\Theta, \Theta\} = 0\) by Proposition \ref{prop:aks}.
- Ghost terms: For \(c^\mu\), \(\{c^\mu \partial_\mu \Phi, \Phi^*\} = \partial_\mu \Phi \partial_\mu \Phi^*\), which cancels with antifield terms via integration by parts.
- Antifield terms: \(\{S^*, c^\mu \partial_\mu S + \eta\} = \partial_\mu S \partial_\mu S^* + \eta S^*\), balanced by ghost contributions.
- Diffeomorphism terms: \(\partial_\nu c^\mu \partial_\mu c^\nu = 0\). Ghost number conservation ensures cancellation \citep{Alexandrov1997}.
\end{proof}

\subsection{Derived Quantization and Boundary Data}

Quantization maps to \(\operatorname{Crit}(S_{\mathrm{BV}})\). Boundary restriction yields:
\[
Z_{\partial \Sigma}[\Gamma] = \int_{\Phi|_{\partial \Sigma} = \Gamma} e^{\frac{i}{\hbar} S_{\mathrm{AKSZ}}}.
\]

\subsection{Interpretation}

The AKSZ–BV framework unifies unitarity and entropy production, with \(\omega_0\) conserving probability and \(dS\) measuring dissipation.

\section{BV–BRST Structure of the RSVP–AKSZ Sigma Model}
\label{app:bv-brst}

\subsection{Field Content and Degrees}

\[
\begin{array}{rclcl}
\Phi &:& \Sigma[0] &\to& \text{Scalar entropy density} \\
\mathbf{v}_i &:& \Sigma[0] &\to& \text{Vector flow field} \\
S &:& \Sigma[0] &\to& \text{Entropy potential} \\
c^\mu &:& \Sigma[1] &\to& \text{Diffeomorphism ghost} \\
\eta &:& \Sigma[1] &\to& \text{Entropy–gauge ghost}
\end{array}
\]
Ghost numbers: \(\operatorname{gh}(c^\mu) = 1\), \(\operatorname{gh}(\eta) = 1\), \(\Phi, \mathbf{v}, S\) at 0. Antifields \(\Psi_A^*\) have degree \(-1 - \operatorname{gh}(\Psi^A)\).

\subsection{Classical Action in BV Form}

\[
S_{\mathrm{BV}} = \int_\Sigma \left[ \Phi (\partial_t S + \mathbf{v} \cdot \nabla S) - D |\nabla S|^2 + \Phi^* (c^\mu \partial_\mu \Phi) + S^* (c^\mu \partial_\mu S + \eta) + v_i^* (c^\mu \partial_\mu v^i - v^j \partial_j c^i) + \frac{1}{2} c_\mu^* c^\nu \partial_\nu c^\mu + \eta^* (c^\mu \partial_\mu \eta) \right].
\]

\subsection{BRST Transformations}

\[
s \Phi = c^\mu \partial_\mu \Phi, \quad s S = c^\mu \partial_\mu S + \eta, \quad s v^i = c^\mu \partial_\mu v^i - v^j \partial_j c^i, \quad s c^\mu = c^\nu \partial_\nu c^\mu, \quad s \eta = c^\mu \partial_\mu \eta.
\]
These satisfy \(s^2 = 0\) up to equations of motion.

\subsection{Classical Master Equation}

See Theorem \ref{thm:bv-master}.

\subsection{Observables and Cohomology}

Observables like \(\mathcal{O}_f = \int_\Sigma f(\Phi, S) d^3x\) lie in BV cohomology, recovering unistochastic statistics.

\subsection{Interpretation}

The BV–BRST formalism generalizes probability conservation. The master equation ensures that stochastic transitions preserve total probability, while entropy production (\(dS\)) measures deviations from unitarity, akin to gauge invariance in field theory \citep{Alexandrov1997}. For non-specialists, this can be understood as a symmetry principle ensuring that probability flows remain consistent across scales, with ghosts and antifields enforcing conservation laws in the presence of stochastic fluctuations.

\section{Derived Quantization and Emergent Probability Law}
\label{app:quantization}

\subsection{BV Path Integral and Gauge Fixing}

\[
Z = \int [\mathcal{D}\Psi][\mathcal{D}\Psi^*] \exp\left( \frac{i}{\hbar} S_{\mathrm{BV}} \right).
\]
Gauge fixing with \(\Psi_{\mathrm{GF}} = \int_\Sigma \bar{c}_\mu G^\mu + \bar{\eta} G_\eta\), where \(G^\mu = \partial_\nu v^\nu\), \(G_\eta = \partial_t S\), yields a BRST-invariant action.

\begin{proof}
The gauge-fixing fermion \(\Psi_{\mathrm{GF}}\) shifts the Lagrangian by \(\delta \Psi_{\mathrm{GF}}\), ensuring BRST invariance. The Faddeev–Popov determinant \(\det(\Delta_{\mathrm{FP}})\) arises from ghost integration \citep{Ikeda1994}.
\end{proof}

\subsection{Integration over Ghosts and Antifields}

Gaussian integration over ghosts yields:
\[
Z = \int [\mathcal{D}\Phi][\mathcal{D}\mathbf{v}][\mathcal{D}S] \exp\left( \frac{i}{\hbar} \int_\Sigma \mathcal{L}_{\mathrm{RSVP}} \right) \times \det(\Delta_{\mathrm{FP}}).
\]

\subsection{Emergent Unistochastic Kernel}

Discretizing \(\Sigma\):
\[
Z \simeq \sum_{\{\Phi_i(t)\}} \exp\left[ \frac{i}{\hbar} \sum_{t,i,j} \Phi_i(t) \Gamma_{ij} S_j(t + \Delta t) \right],
\]
with:
\[
\Gamma_{ij} = \frac{1}{\mathcal{N}} \left| \exp\left( i \Delta t \mathbf{v} \cdot \nabla S - D \Delta t |\nabla S|^2 \right) \right|^2.
\]

\subsection{Born Rule as a Limit Theorem}

See Theorem \ref{thm:bornlimit}.

\subsection{Derived Stack Quantization}

The critical locus \(\operatorname{Crit}(S_{\mathrm{BV}})\) defines a \((-1)\)-shifted symplectic stack, with \(\Gamma_{ij}\) as boundary morphisms.

\subsection{Interpretation}

The BV quantization bridges microscopic unitarity, mesoscopic stochasticity, and macroscopic entropic flow, with the Born rule as a statistical limit.

\section{Semiclassical Expansion and Entropy–Curvature Anomalies}
\label{app:semiclassical}

\subsection{Semiclassical Expansion}

Expanding around \(\Psi_0\):
\[
\Gamma^{(1)} = S_{\mathrm{RSVP}}[\Psi_0] + \frac{i \hbar}{2} \ln \det \mathcal{K}_{AB} + O(\hbar^2).
\]
The fluctuation operator is:
\[
\mathcal{K}_{AB} = \frac{\delta^2 S_{\mathrm{RSVP}}}{\delta \Psi^A \delta \Psi^B}, \quad \mathcal{K}_{\Phi S} = \partial_t S + \mathbf{v} \cdot \nabla S, \quad \mathcal{K}_{S S} = -D \nabla^2 \Phi.
\]

\begin{proof}
Compute second variations:
\[
\frac{\delta^2 S}{\delta \Phi \delta S} = \partial_t S + \mathbf{v} \cdot \nabla S, \quad \frac{\delta^2 S}{\delta S \delta S} = -D \nabla^2 \Phi.
\]
The determinant is regularized via \(\zeta\)-function: \(\ln \det \mathcal{K} = -\zeta'(0)\), with curvature terms from \([\nabla_\mu, \nabla_\nu] S\) \citep{Pantev2013}.
\end{proof}

\subsection{Curvature Coupling and Entropic Torsion}

\[
\Gamma_{\mathrm{anom}} = \frac{\hbar}{24 \pi^2} \int_\Sigma \left( \alpha_1 R \Phi_0 + \alpha_2 T_{\mu\nu}^{\rho} \mathbf{v}_0^\mu \nabla_\rho S_0 + \alpha_3 |\nabla S_0|^2 R \right) d^4x.
\]

\subsection{Effective Unistochastic Kernel}

\[
\Gamma_{ij}^{(\mathrm{eff})} = |U_{ij}|^2 \exp\left[ -\hbar \mathcal{A}_{ij}(\Phi_0, S_0, R, T) \right].
\]

\subsection{Shifted Symplectic Anomaly Form}

\[
d \omega_0 = \hbar \Omega^{(1)}, \quad \Omega^{(1)} = \operatorname{Tr} \left( \mathcal{R} \wedge \mathcal{R} + \mathcal{T} \wedge \mathcal{T} \right).
\]

\subsection{Physical Interpretation}

Anomalies introduce entropy leakage, affecting quantum statistics in curved backgrounds.

\subsection{Summary}

The expansion unifies classical entropy with quantum corrections.

\section{Renormalization–Group Flow of Entropy and Curvature}
\label{app:rg}

\subsection{Scale Dependence}

\[
\mu \frac{d \Gamma}{d \mu} = \left( \beta_D \frac{\partial}{\partial D} + \beta_{\alpha_i} \frac{\partial}{\partial \alpha_i} - \gamma_\Phi \int \Phi \frac{\delta}{\delta \Phi} - \gamma_S \int S \frac{\delta}{\delta S} \right) \Gamma.
\]

\subsection{One–Loop Beta Functions}

\[
\beta_D = \frac{\hbar}{6 \pi^2} D^2 \left(1 - \frac{5}{12} \epsilon \right) - \alpha_3 R D, \quad \beta_{\alpha_1} = \frac{\hbar}{24 \pi^2} (1 - 6D), \quad \beta_{\alpha_2} = \frac{\hbar}{8 \pi^2} D, \quad \beta_{\alpha_3} = \frac{\hbar}{12 \pi^2} (1 - 3D).
\]

\begin{proof}
Using dimensional regularization (\(d = 4 - \epsilon\)), compute one-loop diagrams for \(\Phi, S\) propagators. For \(\beta_D\), the vertex \(\Phi \mathbf{v} \cdot \nabla S\) yields a logarithmic divergence:
\[
\beta_D = \frac{\hbar}{6 \pi^2} D^2 \int \frac{d^d k}{(2\pi)^d} \frac{1}{k^2},
\]
regularized to \(1 - \frac{5}{12} \epsilon\). Curvature terms couple via \(\alpha_3 R\), stabilizing the flow \citep{Wilson1983}. For \(\beta_{\alpha_i}\), loop corrections from curvature terms yield the given coefficients.
\end{proof}

\subsection{Fixed Points}

- **Microscopic**: \(D \to 0\), unitary.
- **Mesoscopic**: \(D_* = 6 \pi^2 \hbar^{-1} \alpha_3 R\), RSVP.
- **Macroscopic**: \(D\) grows, Expyrotic.

\begin{proof}
Solve \(\beta_D = 0\):
\[
\frac{\hbar}{6 \pi^2} D^2 \left(1 - \frac{5}{12} \epsilon \right) - \alpha_3 R D = 0.
\]
The stable fixed point is \(D_* = 6 \pi^2 \hbar^{-1} \alpha_3 R\). The Jacobian \(\partial \beta_D / \partial D\) has negative eigenvalues at \(D_*\), ensuring stability \citep{Wilson1983}.
\end{proof}

\subsection{Flow of Entropy–Curvature Coupling}

\[
\mu \frac{d \Lambda_{\mathrm{ent}}}{d \mu} = (\beta_{\alpha_1} R + \beta_{\alpha_3} |\nabla S|^2) + (\alpha_1 + \alpha_3) \beta_R.
\]

\subsection{Geometric Interpretation}

The RG flow traces a homotopy on \(\omega_0\), connecting unitarity to smoothness.

\subsection{Physical Picture}

\(D\) and \(\alpha_i\) drive the plenum from stochastic unitarity to entropic smoothness.

\subsection{Summary}

The RG flow unifies quantum, thermodynamic, and cosmological scales.

\section{Phase–Space Portrait and Entropic Gradient Flow}
\label{app:phase-space}

\subsection{Entropic Potential}

The RG flow is a gradient descent on:
\[
\mathcal{S}[\Phi, S] = \int_\Sigma \left( \Phi S - \frac{1}{2} D |\nabla S|^2 + \frac{1}{2} \alpha_1 R S^2 + \frac{1}{2} \alpha_3 R |\nabla S|^2 \right) d^4x.
\]

\subsection{Phase–Space Variables}

Coordinates \((q^1, q^2, q^3) = (\Phi, S, D)\), momenta \((p_1, p_2, p_3) = (\nabla \cdot \mathbf{v}, \partial_t S, \partial_t D)\), with symplectic form \(\Omega = \sum_i dq^i \wedge dp_i\).

\subsection{Entropic Phase Portrait}

\[
\frac{d}{d\tau} \begin{pmatrix} \delta D \\ \delta \alpha_i \end{pmatrix} = J \begin{pmatrix} \delta D \\ \delta \alpha_i \end{pmatrix}.
\]
Fixed points are UV-attractive (unistochastic), marginal (RSVP), and IR-attractive (Expyrotic).

\begin{proof}
The Jacobian \(J = \partial \beta / \partial (D, \alpha_i)\) has eigenvalues determined by \(\beta_D, \beta_{\alpha_i}\). At \(D_* = 6 \pi^2 \hbar^{-1} \alpha_3 R\), negative eigenvalues ensure stability \citep{Wilson1983}.
\end{proof}

\subsection{Derived Gradient Flow}

\[
\iota_{Q_{\mathrm{RG}}} \omega_{\mathrm{RG}} = d \mathcal{S}.
\]

\subsection{Cosmological and Informational Outlook}

The flow defines a thermodynamic arrow, with \(D\) as the learning rate.

\subsection{Summary}

The portrait unifies regimes via entropic gradient descent.

\section{Categorical Synthesis of the RSVP–UQM Hierarchy}
\label{app:categorical}

\subsection{Overview}

The categorical hierarchy is:
\[
\begin{tikzcd}
\mathbf{UniStoch} \arrow[r, shift left=0.7ex, "\mathfrak{C}_1"] & \mathbf{RSVPField} \arrow[r, shift left=0.7ex, "\mathfrak{C}_2"] \arrow[l, shift left=0.7ex, "\mathfrak{R}_1"] & \mathbf{BVQuant} \arrow[r, shift left=0.7ex, "\mathfrak{C}_3"] \arrow[l, shift left=0.7ex, "\mathfrak{R}_2"] & \mathbf{RGFlow} \arrow[l, shift left=0.7ex, "\mathfrak{R}_3"]
\end{tikzcd}
\]

\subsection{Layered Structure}

\begin{center}
\begin{tabular}{@{}llll@{}}
\toprule
\textbf{Level} & \textbf{Mathematical Object} & \textbf{Structure Preserved} & \textbf{Physical Interpretation} \\
\midrule
\textbf{I. Unistochastic} & \(\Gamma_{ij} = |U_{ij}|^2\) & Doubly-stochastic & Stochastic microdynamics \\
\textbf{II. RSVP Field} & \((\Phi, \mathbf{v}, S)\) & \(\nabla \cdot \mathbf{v} = 0\) & Entropic continuum \\
\textbf{III. BV–AKSZ} & \(T^*[-1] \mathcal{F}_{\mathrm{RSVP}}\) & 0-shifted symplectic & Quantized entropy \\
\textbf{IV. RG Flow} & \((\mathcal{S}, Q_{\mathrm{RG}})\) & Gradient descent & Scale-dependent smoothing \\
\textbf{V. Global Stack} & \(\mathcal{M}_{\mathrm{RSVP}}\) & Symplectic homotopy & Unified entropy–geometry \\
\bottomrule
\end{tabular}
\end{center}

\subsection{Derived–Stack Composition}

\[
\mathfrak{Q}_{\mathrm{RSVP}} = \mathfrak{C}_3 \circ \mathfrak{C}_2 \circ \mathfrak{C}_1.
\]

\subsection{Homotopy–Commutative Cube}

\[
\begin{tikzcd}
 &[-2em] \mathcal{F}_{\mathrm{RSVP}} \arrow[dr, "\text{Quantize}"] \arrow[dl, "\text{Coarse-grain}"] \\
\mathcal{U}_{\mathrm{unistoch}} \arrow[rr, "\text{Derived BV}"] & & \mathcal{Q}_{\mathrm{BV/AKSZ}} \arrow[dl, "\text{RG flow}"] \\
 & \mathcal{R}_{\mathrm{RG}} \arrow[ul, "\text{Refine}"] &
\end{tikzcd}
\]

\begin{proof}
The cube commutes up to homotopy, with natural transformations \(\eta_i, \epsilon_i\) satisfying \(d \eta_i = d \epsilon_i = dS\) \citep{Pantev2013}.
\end{proof}

\subsection{Monoidal and Probabilistic Structure}

Monoidal products compose subsystems; \(\mathbb{P}: \mathbf{UniStoch} \to \mathbf{Set}\) is lax monoidal.

\subsection{Summary}

The hierarchy forms a derived adjoint tower, unifying stochastic microprocesses with macroscopic geometry.

\section{Categorical Formulation of the RSVP–UQM Tower}
\label{app:categorical-tower}

\subsection{Categories and Objects}

- \(\mathbf{UniStoch}\): Kernels \(\Gamma_{ij}\), stochastic intertwiners.
- \(\mathbf{RSVPField}\): \((\Phi, \mathbf{v}, S)\), entropy-preserving diffeomorphisms.
- \(\mathbf{BVQuant}\): \((\mathcal{F}, \omega_0, \Theta)\), symplectic quasi-isomorphisms.
- \(\mathbf{RGFlow}\): \((\mathcal{S}, Q_{\mathrm{RG}})\), scale-preserving homotopies.

\subsection{Functorial Relationships}

\(\mathfrak{C}_i\) coarse-grain, \(\mathfrak{R}_i\) refine, forming homotopy-adjunctions.

\subsection{Natural Transformations}

\(\eta_i, \epsilon_i\) satisfy \(d \eta_i = d \epsilon_i = dS\).

\subsection{Global Derived Stack}

\[
\mathcal{M}_{\mathrm{RSVP}} = \operatorname{hocolim} \mathcal{D}.
\]

\subsection{Monoidal and Probabilistic Structure}

Monoidal products ensure consistency across scales.

\subsection{Final Summary}

The RSVP–UQM theory unifies stochasticity and geometry via a categorical tower.

\section{Entropy as the Geometry of Indivisibility}

\subsection{Thermodynamic Completion of Quantum Mechanics}

RSVP embeds quantum amplitudes in a geometric medium, resolving measurement issues locally, unlike Bohmian non-locality. The Born rule emerges as an ergodic limit, not a postulate.

\subsection{Relation to Statistical and Information Theories}

RSVP aligns with Jaynes’s maximum entropy principle \citep{Jaynes1957}, minimizing relative entropy:
\[
D_{\text{KL}}(p || q) = \int p \ln \frac{p}{q} \, d^3x.
\]
The action \(\mathcal{S}\) approximates the Kullback–Leibler divergence, with \(D\) as the compression scale \citep{Grunwald2007}.

\subsection{Comparison with Other Interpretations}

\begin{center}
\begin{tabular}{@{}lcccc@{}}
\toprule
\textbf{Feature} & \textbf{ISQM} & \textbf{RSVP} & \textbf{Bohm} & \textbf{GRW} \\
\midrule
Ontology & Stochastic & Entropic field & Deterministic & Stochastic collapse \\
Non-locality & No & No & Yes & No \\
Measurement & Intrinsic jumps & Entropic smoothing & Pilot wave & Collapse terms \\
Born Rule & Emergent & Ergodic limit & Derived & Postulated \\
Empirical Tests & Quantum optics & CMB anisotropies & Bell tests & Collapse rates \\
\bottomrule
\end{tabular}
\end{center}

\begin{proof}
For ISQM, see Theorem \ref{thm:locality}. For RSVP, curvature corrections predict decoherence rates testable in quantum optics. Bohmian non-locality fails Bell tests without additional assumptions. GRW’s collapse rates are experimentally constrained \citep{Ghirardi1986}.
\end{proof}

\subsection{Numerical Methods for RSVP Simulation}

Discretize \(\Sigma\) into a lattice, solve \eqref{eq:PDE1}--\eqref{eq:PDE2} using finite difference methods:
\[
\Phi_i^{n+1} = \Phi_i^n - \Delta t \sum_j J_{ij} + \Delta t \sigma(S_i^n).
\]
Compute \(\Gamma_{ij}\) and verify unistochasticity.

\subsection{Observable Predictions}

RSVP predicts curvature-dependent decoherence rates, testable in quantum optics (e.g., photon interference with gravitational gradients). Cosmological implications include CMB anisotropy spectra modified by \(\Lambda_{\mathrm{ent}}\).

\section{Double-Slit Experiment in ISQM and RSVP}

\subsection{Setup}

Consider a double-slit experiment with slits at \(x_1, x_2\), and a detector screen. ISQM assigns probabilities \(\Gamma_{ij}(t | h)\) based on slit history. RSVP models the slits as cells \(\Omega_1, \Omega_2\), with \(\Phi\) peaked at slit positions.

\subsection{ISQM Prediction}

\[
\Gamma_{ij} = |U_{ij}|^2, \quad U_{ij} = \frac{1}{\sqrt{2}} \begin{pmatrix} e^{i \theta_1} & e^{i \theta_2} \\ e^{i \theta_2} & e^{i \theta_1} \end{pmatrix},
\]
producing interference
